\section{Actual arbitrary-set functional ideals in the same genuine manifold chart}

Let $X$ be a genuine complex manifold of dimension $n$, let $x\in X$,
and let $\Phi$ be supplied genuine local holomorphic coordinates after
the preferred chart $e$ at $x$, with $\Phi(e(x))=0$. Retain the literal
maps $r=\Phi\circ e$ and $q=e^{-1}\circ\Phi^{-1}$, each defined as an
original function. Their inverse identities and holomorphy are used only
on their actual open domains. No smoothness or analyticity is required
of the original subset $N\subset X$.

\begin{lemma}[Entire original functional families in the same supplied chart]
\label{hir:actual-singular-chart-entire-functional-span}
\lean{Hironaka.HironakaActualSingularChartTransport.isInSpanAt_inverseCoordinates_iff}
\uses{trace:actual-manifold-finite-span-adapted-covariance}
\leanok
For any original family $\mathcal S$ of functions on $X$, finite
functional span membership of $F$ at $x$ is equivalent to finite
functional span membership of $F\circ q$ at zero in the entire literal
family $\{s\circ q:s\in\mathcal S\}$. This allows infinite and empty
original families, with no finite-generation or coherence certificate.
\end{lemma}
\begin{proof}
An actual span relation chooses a finite list of original generators.
Regard the same relation as membership in the range span of exactly that
finite list. The genuine finite adapted-chart transport sends the same
original functions to their literal compositions with $q$. Each belongs
to the asserted entire composed family, so generator monotonicity gives
one implication. Conversely, a finite relation in that entire composed
family chooses finitely many composed functions. Choose an original
preimage function for each of these finitely many entries. The literal
function equalities identify the relation with the finite span of those
same preimages composed with $q$. The reverse genuine finite adapted-chart
transport returns the original finite relation, and all chosen original
preimages belong to $\mathcal S$. Generator monotonicity proves the
reverse implication. No infinite sum or global extension is involved.
\end{proof}

\begin{lemma}[Actual vanishing membership for an arbitrary original set in the same inverse chart]
\label{hir:actual-singular-chart-vanishing-membership}
\lean{Hironaka.HironakaActualSingularChartTransport.holomorphicAt_inverseCoordinates_iff,Hironaka.HironakaActualSingularChartTransport.vanishingIdeal_inverseCoordinates_iff}
\uses{hir:actual-singular-chart-entire-functional-span}
\leanok
For every original subset $N\subset X$ and function $F:X\to\mathbb C$,
\[
 F\in(\mathcal I_N)_x\quad\Longleftrightarrow\quad
 F\circ q\in(\mathcal I_{q^{-1}N})_0.
\]
The set on the right is the actual entire inverse image under the literal
original inverse coordinate function, rather than an assumed analytic
presentation. In particular holomorphy of $F$ at $x$ and of $F\circ q$
at zero are equivalent.
\end{lemma}
\begin{proof}
The genuine finite adapted-chart span theorem applied to the one-element
family containing $F$ gives exactly the asserted equivalence of
holomorphy. The actual inverse chart $q$ tends to $x$ at zero by
continuity and its real local inverse identity; hence a genuine
neighbourhood on which $F$ vanishes on $N$ pulls back to a genuine
neighbourhood on which $F\circ q$ vanishes on $q^{-1}N$.
Conversely the actual forward map $r$ tends to zero at $x$.
Pull a local zero relation for $F\circ q$ back through $r$ and intersect
with the true open coordinate domain of $x$. On this domain the real
identity $q(r(z))=z$ holds. It converts both the set membership and the
function value into the original membership in $N$ and the original
value $F(z)$. Together with the holomorphy equivalence this proves the
reverse original vanishing-ideal implication.
\end{proof}

\begin{lemma}[Every model vanishing function has its actual original representative in the same chart]
\label{hir:actual-singular-chart-original-function-representative}
\lean{Hironaka.HironakaActualSingularChartTransport.original_representative_of_model_vanishing}
\uses{hir:actual-singular-chart-vanishing-membership}
\leanok
Suppose the source of $\Phi$ lies in the actual target of the preferred
chart $e$. If an original model function $g$ belongs to
$(\mathcal I_{q^{-1}N})_0$, the literal original function $a=g\circ r$
belongs to $(\mathcal I_N)_x$, and $a\circ q=g$ on a genuine
neighbourhood of zero. No original representative is supplied as a
certificate or extended holomorphically to the entire manifold.
\end{lemma}
\begin{proof}
The true target of $\Phi$ is an open neighbourhood of zero. For each
point $u$ in it, the actual source inclusion ensures that
$e^{-1}(\Phi^{-1}(u))$ belongs to the true original coordinate domain,
and the literal right inverse identity gives $r(q(u))=u$.
Thus $(g\circ r)\circ q=g$ on this genuine neighbourhood.
Original vanishing membership is invariant under equality on a
neighbourhood: holomorphy is invariant under that equality, and the
local zero relation is simply restricted to its intersection.
Therefore $(g\circ r)\circ q$ has the stated model vanishing membership.
The preceding inverse-chart theorem returns actual membership of the
same literal $g\circ r$ in $(\mathcal I_N)_x$. Its arbitrary values
outside the coordinate neighbourhood play no role; no global
holomorphic extension has been asserted.
\end{proof}

These dictionaries retain a supplied genuine chart and arbitrary original
sets and functions. Constructing the product commutation for a genuine
submersion and then deriving its whole singular vanishing-ideal pullback
is a further geometric application, not an assumed ideal equality.
